\section*{Overview}

2-SAT is the form of Boolean satisfiability that still behaves like graph theory. Once every clause has exactly two
literals, the problem can be turned into an implication graph and solved with SCC decomposition in linear time.

\section*{When to Use It}

Use 2-SAT when:

\begin{itemize}
  \item choices are binary,
  \item constraints look like ``at least one of these two literals must hold'',
  \item the problem is really about compatible or incompatible pairs of decisions.
\end{itemize}

This appears in switch toggles, orientation choices, scheduling with binary states, and selection problems with pairwise
logic constraints.

\section*{Core Idea}

Represent each variable \(x\) by two nodes:
\[
x \quad\text{and}\quad \lnot x.
\]
A clause \((a \lor b)\) is equivalent to the two implications:
\[
\lnot a \Rightarrow b,
\qquad
\lnot b \Rightarrow a.
\]
Add those edges to the implication graph, then compute SCCs.

\section*{Key Insight}

If \(x\) and \(\lnot x\) lie in the same SCC, then each implies the other, so the formula forces both true and false at
once. Otherwise, the SCC topological order gives a satisfying assignment.

\section*{Operations / Main Technique}

\begin{itemize}
  \item encode each literal as one graph node,
  \item add implications for every clause,
  \item compute SCCs,
  \item assign each variable by SCC order.
\end{itemize}

\paragraph{Worked example.}
The clause \((x \lor \lnot y)\) produces:
\[
\lnot x \Rightarrow \lnot y,
\qquad
y \Rightarrow x.
\]
These are the two graph edges the solver actually needs.

\section*{Correctness Intuition}

Every clause \((a \lor b)\) fails only when both \(a\) and \(b\) are false, so the implication form precisely captures
which literals must become true when another literal is false. SCCs detect contradictions, and the reverse topological
order ensures that once one side of a variable pair is chosen, all of its implications are respected.

\section*{Complexity Analysis}

For \(n\) variables and \(m\) clauses:
\begin{itemize}
  \item graph size: \(2n\) vertices and \(2m\) edges,
  \item total time: \(O(n + m)\),
  \item memory: \(O(n + m)\).
\end{itemize}

\section*{Implementation}

The code includes a compact solver with:

\begin{itemize}
  \item clause insertion,
  \item implication insertion,
  \item satisfiability check,
  \item one recovered assignment.
\end{itemize}

\section*{Common Pitfalls}

\begin{itemize}
  \item Mixing variable indices with literal indices incorrectly.
  \item Forgetting that the assignment is read from SCC order, not from raw DFS order.
  \item Trying to force constraints that are not reducible to binary clauses.
\end{itemize}

\section*{Variants / Extensions}

\begin{itemize}
  \item Encode ``at most one'' constraints with several binary clauses in small cases.
  \item Use the implication graph to explain forced assignments and contradictions.
  \item General SAT, which is far harder and not covered by this graph trick.
\end{itemize}

\section*{Practice Problems}

\begin{itemize}
  \item Assign binary states under pairwise logical constraints.
  \item Choose orientations or colors from two options with incompatibility rules.
  \item Giant Pizza and similar implication-graph tasks.
\end{itemize}

\section*{References}

\begin{itemize}
  \item \href{https://cp-algorithms.com/graph/2SAT.html}{cp-algorithms: 2-SAT}.
  \item \href{https://usaco.guide/adv/SCC}{USACO Guide: SCC and 2-SAT}.
  \item \href{https://codeforces.com/catalog?locale=en}{Codeforces Catalog}.
\end{itemize}
